\begin{partbacktext}
\part{Preliminaries}
\noindent
This part sets up the analytic and geometric language used throughout the book. It recalls the basic theory of quasi-plurisubharmonic functions, non-pluripolar products, envelope operators, geodesics of singular potentials and the toric-convex dictionary in the ample case. A reader already familiar with pluripotential theory on compact K\"ahler manifolds may skim parts of \cref{chap:pshfunctions,chap:npp}, but the later chapters use the precise conventions fixed here.

\cref{chap:pshfunctions} recalls plurisubharmonic and quasi-plurisubharmonic functions, pluripolar sets, the plurifine topology, Lelong numbers, multiplier ideals, analytic singularities, positivity of currents and singular Hermitian metrics. The presentation is concise, but we include proofs of several points whose conventions are important later, especially those involving Lelong numbers, multiplier ideals and quasi-equisingular approximation.

\cref{chap:npp} introduces the non-pluripolar product and the basic tools for manipulating Monge--Amp\`ere masses of singular currents.

\cref{chap:envelopes} develops the two envelope operators that organize the singularity theory of the book. The $P$-envelope records the singularity type relevant to non-pluripolar mass, while the $\mathcal I$-envelope records the singularity type detected by multiplier ideals. Their interaction is one of the main themes of the book.

\cref{chap:rays} studies subgeodesics, geodesics and geodesic rays in spaces of singular potentials. It also introduces the finite-energy space $\mathcal E^1(X,\theta;\phi)$ and the metric structures that will later be used to define convergence of singularity types.

Finally, \cref{chap:toric_ample} treats the ample toric case. There the pluripotential theory of toric metrics is translated into convex analysis on the associated polytope. This chapter is logically independent of much of the rest of the book, but it provides a model for the big toric theory developed later in \cref{chap:toric_big}.
\end{partbacktext}
